In this first part we discuss the basics of language modeling and machine learning. We start with basic principles of how language can be modeled, using tokens and embeddings. Next, the attention concept is discussed, a concept that is at the foundation of the transformer architecture, and the success of large language models. We will also review two important machine learning techniques: supervised learning and reinforcement learning. In the next chapter, we become more concrete, and discuss the steps of a typical training pipeline, from pretraining and finetuning to inference. We also discuss common benchmarks in that chapter. The final content chapter of this first part is about the {\em large} part of large language models: we discuss how performance of LLMs scales with parameter count, training data and compute, and we delve into efficient models that benefit from scale while minimizing cost. At the end of the part, we  turn to a hands-on  case study, where you will learn how everything comes together when you create your own tiny large language model from scratch. 
